\begin{abstract}
大语言模型（LLM）擅长读懂长文档、组织论证与写作，却在投研场景中有三类不可接受的失败：编造或记错财务数字、在增速与估值上心算出错、凭参数记忆补写没有出处的外部信息。改写环节还会引入新的数字错误，而研报质量本身缺乏可复核的评价，使“改进”难以度量。本文完整介绍 \sysname{}（\textbf{L}LM-driven \textbf{E}quity Research with \textbf{D}eterministic \textbf{G}rounding, \textbf{E}valuation and \textbf{R}eporting）——一个面向 A 股上市公司的可追溯投研多智能体系统。它围绕财报分析、自动化估值建模与研究报告质量评估三个主选题构建，并覆盖数据提取、产业链分析与研究报告纠错核查三个子选题；它能把一份年报加工为 25～32 页、60 个以上图表的卖方体例深度研报与公式联动的估值工作簿。

系统的核心原则是\textbf{理解与计算分离}：模型负责读、抄、写、判断，确定性程序负责校验、计算、拦截。我们把这一原则落实为一组可测试的不变量与对应机制：（1）\textbf{证据约束的双通道抽取}——模型抄录的每个数字必须在所写页码的原文中按数字边界找到，规则通道独立取数交叉校验，冲突不静默取舍而是阻断估值；（2）\textbf{单一真值表}驱动的分析与估值链——分部量价预测、三表联动、三情景 FCFF 折现、逐项股权价值桥、反向 DCF 与可比公司，评级由程序给出；（3）按研报\textbf{展示精度}（半个末位）设定容差的纠错与护栏，误报控制优先于召回；（4）\textbf{数字锁}，保证模型改写只能删减、不能新增或改动数字；（5）由确定性编排器驱动、五个上下文隔离的无头 Claude Code 会话（研究助理、分析师、主笔、复核编辑、排版设计师）与 \textbf{11 道程序闸门}构成的研报流水线，先搜集外部资料、后分析写作；（6）四层评分标准与\textbf{对生成方不可见}的盲评闭环，并用重复评审、程序制造缺陷与人工复核三组证据\textbf{校准评审模型本身}。

在三家公司的真实年报上，模型抄录并通过证据校验的 301 个数值中，规则通道独立取到且一致的有 280 个（93.0\%），冲突为 0，勾稽 56/56 全部通过；流水线引入先搜集与深度专题后，两家公司深度研报的质检得分由 74、67 分提升至 90、88 分（同一评分标准、单次盲评）；评审模型重复盲评的逐条一致率为 86\%～91\%，判出问题经人工确认的精确率为 84\%。本文逐一拆解每个模块的动机、形式化定义、算法、实现与设计权衡，讨论系统的演进史与失败的尝试，并如实报告证据的规模与局限。所有程序计算均可在保存的输入上离线逐位复算。
\end{abstract}
